\caption{Examples of horizontal grid geometries and variable staggering used
in atmospheric models. Panels (a--c) show regular latitude--longitude,
equiangular gnomonic cubed-sphere, and icosahedral-dual meshes in the same
orthographic view. The dual uses native OASIS polygon corners on a coarse
mesh without spring relaxation; heavier outlines identify visible pentagons
(twelve globally). Mesh densities are chosen for legibility and do not imply
equal effective resolution. Panels (d--f) independently illustrate A-, C-,
and Z-grid arrangements on the same local Cartesian mesh. Here $h$ denotes a
generic mass/thickness-like scalar, $u$ and $v$ are horizontal velocity
components, and $\delta$ and $\zeta$ are horizontal divergence and vertical
relative vorticity. Dots identify stored scalars, and arrows identify velocity
locations and component directions, not a wind field. C-grid velocity
components are normal to their respective cell faces; Z-grid velocity is
reconstructed from divergence and vorticity. The two groups imply no pairings
between geometry and staggering.}
\label{fig:horizontal-grid-geometry-staggering}
